%\documentclass[a4paper,twoside,12pt]{book}
\documentclass[twoside, openright, 11pt]{report}
\usepackage[spanish]{babel}
\usepackage[utf8]{inputenc}
\usepackage{graphicx} %Para poder insertar las imagenes
\usepackage{enumitem}%Para trabajar con los enumerates
\usepackage[table,xcdraw]{xcolor}%Tablas autogeneradas con tablesgenerator
\usepackage{csquotes}%Para poder entre comillas los textos
\usepackage{wrapfig}%Imagenes junto al texto
%\usepackage{url}
%\usepackage{blindtext}

\setquotestyle{english}%Para que el quote me lo haga con `` y no con <<
\begin{document}

\begin{titlepage}

\newlength{\centeroffset}
\setlength{\centeroffset}{-0.5\oddsidemargin}
\addtolength{\centeroffset}{0.5\evensidemargin}
\thispagestyle{empty}

\noindent\hspace*{\centeroffset}\begin{minipage}{\textwidth}

\centering
\includegraphics[width=0.9\textwidth]{imagenes/logo_ugr}\\[1cm]

\textsc{ \Large TRABAJO FIN DE GRADO\\[0.2cm]}
\textsc{ GRADO DE INGENIERÍA EN INFORMÁTICA}\\[1cm]
% Upper part of the page
%
% Title
{\huge\bfseries Desarrollo de sistema basado en geolocalización dirigido a personas con Alzheimer y su entorno\\
}
\noindent\rule[-1ex]{\textwidth}{3pt}\\[3.5ex]
{\large\bfseries Desarrollo de Software}
\end{minipage}

\vspace{0.3cm}
\noindent\hspace*{\centeroffset}\begin{minipage}{\textwidth}
\centering

\textbf{Autor}\\ {ELENA CHAVES HERNÁNDEZ}\\[2.5ex]
\textbf{Directores}\\
{CARLOS CANO GUTIÉRREZ}\\
{CARLOS RODRÍGUEZ DOMÍNGUEZ}\\[2cm]
\includegraphics[width=0.3\textwidth]{imagenes/logo-ceuta.jpg}\\[0.1cm]
%\includegraphics[width=0.15\textwidth]{imagenes/feetce.jpg}\\[0.1cm] %otra imagen que tenía por ahí del logo de la facultad donde se ve mejor, pero no queda igual de bien
\textsc{Facultad de Educación, Economía y Tecnología de Ceuta}\\
\textsc{---}\\
Granada, 15 de Junio de 2021
\end{minipage}
\end{titlepage}
\let\cleardoublepage\clearpage

\chapter*{}
\begin{flushright}
\textit{Dedicado a \\
...}
\end{flushright}
\thispagestyle{empty}

\chapter*{Resumen} % si no queremos que añada la palabra "Capitulo"
\addcontentsline{toc}{chapter}{Resumen} % si queremos que aparezca en el índice
\markboth{RESUMEN}{RESUMEN} % encabezado
\thispagestyle{empty}
En este proyecto se procederá con el diseño y el desarrollo de un sistema basado en geolocalización dirigido a personas con Alzheimer y su entorno... %incluimos archivo que contiene la sección del resumen

\chapter*{Abstract} % si no queremos que añada la palabra "Capitulo"
\addcontentsline{toc}{chapter}{Abstract} % si queremos que aparezca en el índice
\markboth{ABSTRACT}{ABSTRACT} % encabezado
\thispagestyle{empty}
Resumen en inglés...%incluimos el archivo que contiene el resumen en inglés

\tableofcontents % indice de contenidos

\cleardoublepage
\addcontentsline{toc}{chapter}{Lista de figuras} % para que aparezca en el indice de contenidos
\listoffigures % indice de figuras

\cleardoublepage
\addcontentsline{toc}{chapter}{Lista de tablas} % para que aparezca en el indice de contenidos
\listoftables % indice de tablas

\chapter{Introducción}\label{cap.introduccion}

  %No me gusta esto de aquí vacío
 
 %IoT; eHealth; sistema ubicuo.
 
 %\textit{IoT} (\textit{Internet of Things}, en español ``Internet de las Cosas'') Las siglas así conocidas no se ponen en cursiva.
 
 

  \section{Motivación}\label{sec.motivacion}
  
  %Estudiar las problemáticas típicas de las personas con Alzheimer

    La demencia se puede definir como el deterioro adquirido de las capacidades cognitivas que entorpece la realización satisfactoria de las actividades diarias. El principal causante de demencia, que se estima entre un 60 y un 70 por ciento, es la Enfermedad de Alzheimer (EA), una enfermedad degenerativa que se da habitualmente en pacientes mayores de 65 años. Es importante en este tipo de enfermedades considerar la poca información que se tiene sobre la enfermedad y que, hasta ahora, no tiene cura. Esto se traduce en que no sólo afecta al sujeto, sino también a sus cuidadores y familiares por el grado de dependencia y discapacidad que se puede llegar a sufrir durante la EA.\cite{BarreraLopez2018}\\
 
    %	Cambiar lo que sigue, no me gusta 
    %	Es un campo bastante inexplorado dada su complejidad (es un campo desconocido porque es bastante incierto)

    %Esta enfermedad, caracterizada por el cúmulo de ciertas proteínas en el cerebro, condiciona un deterioro cognitivo progresivo e invariable, afectando principalmente a la memoria reciente. 
 En este trabajo, nos centramos en los pacientes que padecen la EA en una fase temprana y su entorno. En esta fase se observa, aparte de la pérdida de memoria reciente, episodios de ansiedad y falta de asociatividad, concentración, sueño y apetito, entre otros. Esto supone tanto para la persona que padece EA como para su entorno un aumento de responsabilidades en el cuidado, dándose una gran pérdida de autonomía en todas las partes implicadas y llegando a generar problemas graves de ansiedad y depresión.\\
 
 
 %importante poner en la motivación lo que puecde ocurrir con la pérdida de memoria, como la deorientación y pérdida de unapersona en la calle
  %La siguiente frase la quiero cambiar
   Una plataforma de servicios orientada y creada especialmente atendiendo a las necesidades de todas las personas implicadas, puede ofrecer no sólo tranquilidad al llevar un seguimiento, monitorización o información relevante sobre la persona al cuidado, sino también una mayor autonomía. Todo esto se traduce en un aumento de la calidad de vida, en concreto, colocar un pequeño dispositivo de localización sobre la persona que padece EA, disminuye la dependencia del sujeto. Utilizando IoT junto a elementos como una etiqueta NFC, puede ayudar a llevar a cabo tareas cotidianas, como olvidar qué comprar, que se empiezan a ver afectadas desde el inicio de la EA.\\

 

  \section{Objetivos}\label{sec.objetivos}
  %sistema de ayuda
  %mayor autonomía
  
  \subsection{Objetivo General}
  
  El objetivo general de este proyecto es crear un sistema completo con geolocalización haciendo uso de tres dispositivos. Estos dispositivos incluyes un pequeño dispositivo hardware IoT especialmente diseñado para el proyecto a partir de componentes disponibles en el mercado, y dos terminales móviles. Con esto se pretende conseguir una mayor autonomía tanto para el paciente con Alzehimer como para sus cuidadores.
  %Ubicuo, eHealth, servicios
  
  \subsection{Objetivos específicos}
  
  Los objetivos específicos de este proyectos son:
  
  \begin{enumerate}[label=$\bullet$ ]
		\item Hacer uso de los componentes hardware orientados al IoT disponibles en el mercado para construir un dispositivo que se adecúe a nuestras necesidades.
		\item Crear una aplicación especialmente diseñada para personas que sufren EA y su entorno.
		\item Unificar las tecnologías actuales para crear un sistema completo.
  		\item Compartir con la comunidad la información, y los avances recogidos en este proyecto.
  \end{enumerate}
  
  
  \section{Estructura de la memoria}\label{sec.estructura de la memoria}
  %Explicar de qué va la memoria capítulo por capítulo brevemente (creo que me he pasao de breve)
  
  En este primer capítulo introductorio se exponen la motivación y los objetivos del proyecto, al igual que se los recursos utilizados y el desarrolo temporal.\\
  
  \par
  En el Capítulo \ref{cap.estado del arte}, se realiza un análisis documental sobre los fundamentos tanto teóricos como técnicos sobre los sistemas y aplicaciones existentes cuyos objetivos se asemejan a los de este proyecto. Se conl¡cluirá con las carencias y las nuevas aportaciones para mejorar la calidad de estos sistemas y aportarle un valor añadido a nuestro proyecto.\\
  
  \par
  En el Capítulo \ref{cap.diseno y descripcion del sistema}, se establecen los requisitos previos al diseño del sistema, así como los requisitos y la arquitectura general de este. También se añade una sección donde se contemplan los distintos Mockups diseñados según las necesidades de los distintos usuarios del sistema.\\
  
  \par
  A partir de lo expuesto en el Capítulo \ref{cap.diseno y descripcion del sistema}, se especifican en el Capítulo \ref{cap.prototipos y desarrollo} los distintos prototipos que se han ido construyendo y desarrollando de manera incremental hasta llegar al prototipo final.\\
  %describir cada prototipo final brevemente para orientarlo?
  \par
  En el Capítulo \ref{cap.conclusiones y mejoras futuras} se hará un análisis sobre las metas y objetivos cumplidos, así como una valoración sobre algunos aspectos y enfoques futuros de cara a mejorar el sistema.\\ 
  
  \par
  El Capítulo \ref{cap.conclusions and future works} tiene el mismo contenido que el Capítulo \ref{cap.conclusiones y mejoras futuras} traducido al inglés.
  

 \section{Recursos utilizados}\label{sec.recursos}
 %Se irá rellenando conforme vayamos utilizando diferente tecnología (describir con un par de líneas)
 %poner en ionic que se utiliza Angular y Node.js
 
 %El sistema operativo es Microsoft Windows 10 Pro Versión 10.0.19041 compilación 19041
 
 \subsection*{Recursos software}
 
 \begin{enumerate}[wide=\parindent, label= ]
 	\item Angular: Framework modular y escalable diseñado para el desarrollo de aplicaciones web utilizando el patrón MVC (Modelo-Vista-Controlador). Su lenguaje de programación es TypeScript, una ampliación de JavaScript. Versión: 8.3.25
 	\item AdobeXD: Herramienta de edición gráfica orientada a la creación de interfaces tanto web como móviles. Se utiliza para la creación de Mockups. Versión: 36.2.32.5
 	\item Node.js: Entorno de ejecución de JavaScript diseñado para la creación de aplicaciones web escalables y orientado a eventos asíncronos. En este proyecto lo utilizaremos tanto como motor JavaScript  para Angular, como para crear servicios REST. Versión: 14.15.4
 	\item MySQL Server: Sistema de gestión de bases de datos relacionales usando el modelo cliente-servidor. Versión: 8.0.23
 	\item MySQL Workbench: IDE para trabajar con el sistema SQL. Versión: 8.0
 	\item JetBrains Product Pack for Students: Paquete gratuito de IDEs ofrecido a estudiantes de mano de JetBrains. Haremos uso de PyCharm y WebStorm. Versiones: 2020.3
 	\item Django: Framework de alto nivel para el desarrollo web utilizando el patrón MVC (modelo-Vista-Controlador). Desarrollado en Python, lo utilizaremos para crear servicios. Versión: 3.1.7
 	\item Python: Lenguaje de programación interpretado. Versión: 3.8.6
 	\item JavaScript: Lenguaje de programación para web.
 	\item MongoDB: Sistema de gestión de bases de datos NoSQL usando el modelo cliente-servidor orientado a documentos de código abierto. Versión: 4.4.4
 	%\item MongoDB Compass: al final no lo he usado na mas que pa bichear más cómodo que con mongo
 	\item Ionic: Framework para el desarrollo de aplicaciones multiplataforma. Como framework de interfaz utilizaremos Angular con motor Node.js. Versión: 6.12.4
 	\item Postman: Herramienta para hacer tests de APIs REST. Versión: 8.1.0
 	\item Mendeley Desktop: Gestor de referencias bibliográficas multiplataforma. Versión 1.19.4
 	\item \LaTeX\ : Sistema de composición de textos. El editor utilizado es TeXworks y la distribución MikTeX.
 	\item git: Sistema de control de versiones distribuido. Los repositorios se localizan en GitHub. Versión: 2.28.0.windows.1
 \end{enumerate}
 
 \subsection*{Recursos hardware}
 
 \begin{enumerate}[label=]
 	\item 
 \end{enumerate}

\section{Planificación temporal}\label{sec.planificacion temporal}
%¿Cronograma? No, gracias. Fases y subfases de trabajo justificadas.


\begin{table}[h!]
\resizebox{\textwidth}{!}{
\begin{tabular}{l|l|l|l|l|l|l|l|l|l|l|l|l|l|l|l|l|l|l|}
\cline{2-19}
                                                                                         & \multicolumn{4}{c|}{\cellcolor[HTML]{2ACB2A}Febrero}                                                      & \multicolumn{4}{c|}{\cellcolor[HTML]{2ACB2A}Marzo}                                                        & \multicolumn{4}{c|}{\cellcolor[HTML]{2ACB2A}Abril}                                                        & \multicolumn{4}{c|}{\cellcolor[HTML]{2ACB2A}Mayo}                                                         & \multicolumn{2}{c|}{\cellcolor[HTML]{2ACB2A}Junio}  \\ \hline
\multicolumn{1}{|l|}{\cellcolor[HTML]{E9E441}Análisis de requisitos}                     & \cellcolor[HTML]{00D2CB} & \cellcolor[HTML]{00D2CB} & \cellcolor[HTML]{00D2CB} &                          &                          &                          &                          &                          &                          &                          &                          &                          &                          &                          &                          &                          &                          &                          \\ \hline
\multicolumn{1}{|l|}{\cellcolor[HTML]{E9E441}Estudio del software y hardware disponible} &                          &                          & \cellcolor[HTML]{00D2CB} & \cellcolor[HTML]{00D2CB} & \cellcolor[HTML]{00D2CB} &                          &                          &                          &                          &                          &                          &                          &                          &                          &                          &                          &                          &                          \\ \hline
\multicolumn{1}{|l|}{\cellcolor[HTML]{E9E441}Diseño}                                     &                          &                          &                          &                          & \cellcolor[HTML]{329A9D} & \cellcolor[HTML]{329A9D} & \cellcolor[HTML]{329A9D} & \cellcolor[HTML]{329A9D} & \cellcolor[HTML]{329A9D} & \cellcolor[HTML]{329A9D} &                          &                          &                          &                          &                          &                          &                          &                          \\ \hline
\multicolumn{1}{|l|}{\cellcolor[HTML]{E9E441}Implementación de prototipos}               &                          &                          &                          &                          &                          &                          &                          &                          &                          & \cellcolor[HTML]{329A9D} & \cellcolor[HTML]{329A9D} & \cellcolor[HTML]{329A9D} & \cellcolor[HTML]{329A9D} & \cellcolor[HTML]{329A9D} & \cellcolor[HTML]{329A9D} & \cellcolor[HTML]{329A9D} &                          &                          \\ \hline
\multicolumn{1}{|l|}{\cellcolor[HTML]{E9E441}Presentación y conclusiones}                &                          &                          &                          &                          &                          &                          &                          &                          &                          &                          &                          &                          &                          &                          &                          & \cellcolor[HTML]{34696D} & \cellcolor[HTML]{34696D} & \cellcolor[HTML]{34696D} \\ \hline
\rowcolor[HTML]{F8A102} 
\multicolumn{1}{|l|}{\cellcolor[HTML]{E9E441}Confección de la memoria}                   &                          &                          &                          &                          &                          &                          &                          &                          &                          &                          &                          &                          &                          &                          &                          &                          &                          &                          \\ \hline
\end{tabular}}
\end{table}
  
%incluimos el archivo donde se encuentra el capítulo de la introducción por puntos

\chapter{Estado del arte}\label{cap.estado del arte}

  \section{Fundamentos}\label{sec.fundamentos}
  
  Desde hace más de una década se utiliza el término eHealth consistente en, según la OMS desde 2005,  \enquote{el apoyo que la utilización costoeficaz y segura de las tecnologías de la información y las comunicaciones ofrece a la salud y a los ámbitos relacionados con ella,  con  inclusión  de  los  servicios  de  atención  de  salud,  la  vigilancia  y  la  documentación  sanitarias,  así como la educación, los conocimientos y las investigaciones en materia de salud}\cite{OrganizacionMundialdelaSalud2005}. Este término encaja perfectamente con nuestro proyecto, una plataforma de servicios ofrecido especialmente para atender la salud de los pacientes objetivos.\\
  
  \par
  El avance en las tecnologías aplicadas a la eHealth nos lleva a otro término muy extendido y que engloba gran parte de los elementos que forman parte de nuestro día a día, el \enquote{IoT}, del inglés, \enquote{\textit{Internet of Things}}. Este término hace referencia a la interconexión existente entre todo tipo de dispositivos, desde un sensor de movimiento o una cámara de vigilancia hasta cualquier tipo de dispositivo móvil complejo o incluso industrial, como puede ser un Smartphone.\\
  
  \par
  Para realizar la conexión entre estos dispositivos que forman parte de una red entre sí y con la información con la que trabajan, están cada vez más extendidos los denominados \enquote{\textit{microservicios}}, que dividen el software en elementos muchos más pequeños, lo que aporta una gran escalabilidad, la posibilidad de trabjar en distintos lenguajes según el servicio que se quiera ofrecer o una mayor capacidad de recuperación, ya que los microservicios pueden trabajar de forma independiente; en contraposición, añade complejidad al diseño, ya que tienen que estar coordinados entre sí y las pruebas pueden llegar a ser algo tediosas.\\
  
  \par
  En cuanto a la tecnología actual sobre computación y almacenamiento de datos, podemos decir que el futuro es lo que conocemos como \enquote{Cloud Computing}. Actualmente, se ofrecen todo tipo de servicios en la nube a nivel computacional; lo más importante es que únicamente se paga por los recursos utilizados y te permite de manera dinámica aumentar o disminuir los servicios contratados en la mayoría de los casos.\\
  
  \par
  
  Por otro lado, el desarrollo software se lleva a cabo actualmente a través de \enquote{Frameworks}, tanto para el desarrollo de aplicaciones como de servicios. Los más utilizados actualmente son \enquote{Ionic} y \enquote{React Native} que permiten la creación de aplicaciones multiplataforma. En cuanto a la prueba de sistemas, cada vez está más extendida la técnica de Sandwich, que supone un híbrido entre un test del sistema completo a cada componente (de arriba hacia abajo) y de cada componente hasta el sistema completo. Esta técnica persigue probar en paralelo la interfaz de usuario y los componentes software del sistema para en una etapa final, hacer pruebas sobre el sistema completo; esto permite ganar tiempo en pruebas, pero necesita bastantes drivers para las pruebas y puede tener un coste bastante elevado.\\
  
  
  
  

\section{Análisis de sistemas existentes}\label{sec.analisis de sistemas existentes}
   %En cada subsección, explicar quién lo propone y por qué lo propone. Cuándo lo propone, sus características. Ventajas e inconvenientes. Hablar de sistemas informáticos, pero tb de sistemas asistenciales.
   
  \subsection*{Jiobit}
  
  	Jiobit es un sistema de seguimiento en tiempo real que fue diseñado por Jhon Renaldi, que tras perder a su hijo en un parque en 2015, buscó dispositivos para localizar en tiempo real, pero no encontró ningún dispositivo \enquote{wearable} adecuado. Buscaba que fuera elegante, tuviera una batería duradera y no dependiera de un dispositivo, ya fuera por necesitar de conexión bluetooth o de cualquier otro tipo. \\
  	  	\par
  	\begin{wrapfigure}{l}{0.25\textwidth}
  		\centering
		\includegraphics[width=0.25\textwidth]{imagenes/Jiobit.jpg}
	\end{wrapfigure}
  	

  	Aunque este dispositivo fue diseñado para su uso en niños, por su versatilidad pronto se extendió a mascotas, adultos y personas mayores. Tiene el tamaño de una galleta y dispone de un botón de alerta para avisar en tiempo real, el dispositivo se conecta a un servicio en la nube que a su vez conecta con los dispositivos a los que esté asociado Jiobit. \\
  	\par
  	
  	\textbf{Ventajas:}
  	
\begin{enumerate}[label=-]
  		\item No necesita de otro dispositivo para enviar la posición actual al servidor.
  		\item Es ligero y fácil de colocar.
  		\item Cuenta con un botón de aviso.
  		\item Permite establecer sitios seguros.
\end{enumerate}

  	\textbf{Inconvenientes:}
  	
\begin{enumerate}[label=-]
  		\item Hay que pagar mensualmente por el servicio.
  		\item Es fácil que el paciente se deshaga del dispositivo al ser un \enquote{wearable} bastante visible.
  		\item Puede que se olvide la recarga del dispositivo.
\end{enumerate}

\subsection*{Alexa Eldercare Toolbox\cite{Tan2020}}

	Este sistema fue propuesto por investigadores del \enquote{Instituto Tecnológico de Stevens} en 2020, partiendo de estudios en los que la población anciana prefiere quedarse en su hogar. tomando como base el coste anual que puede suponer el cuidado de una persona mayor, se diseñó un sistema utilizando un dispositivo Amazon Echo Dot. De esta manera, también se evita el cuidado en casa, reduciendo el riesgo de contagio de COVID-19.\\
	
	\par
	La potencia de este sistema es la amplia gama de necesidades que puede cubrir, algunas de las que se cubrieron con este sistema fueron: ayudar en la realización de tareas diarias a través de instrucciones, control del tiempo en la cocina, ayuda con la dieta o riesgos de caídas. \\
	
	\par
	  	\textbf{Ventajas:}
  	
\begin{enumerate}[label=-]
  		\item Cubre muchísimas necesidades dentro del hogar.
  		\item Se puede configurar facilmente a través de una aplicación.
  		\item Totalmente personalizable.
  		\item Contacto directo en tiempo real.
\end{enumerate}

  	\textbf{Inconvenientes:}
  	
\begin{enumerate}[label=-]
  		\item No avisa cuando sale de casa.
  		\item Depende de una conexión a internet.
  		\item Puede que se olvide la recarga del dispositivo.
  		\item Amazon Echo se debe colocar en un sitio estratégico de cada domicilio para que esté en continuo contacto.
  		\item El paciente puede rehusarlo, puesto que es un cambio significativo para la vida cotidiana.
\end{enumerate}

%No sé de qué más tipos de sistemas hablar, hay de todo.

  \section{Análisis de aplicaciones móviles}\label{sec.analisis de aplicaciones moviles}
  %En cada subsección, explicar quién lo propone y por qué lo propone. Cuándo lo propone, sus características. Ventajas e inconvenientes.
  
  \subsection*{Alzheimer Master}
  
  	Esta aplicación, creada por un empresa Húngara y publicada en 2017 está disponible para Android. Su función es crear notificaciones con mensajes de voz o video personalizados y grabas la reacción de la persona afecta ante el mensaje. Su objetivo es combatir la soledad y ofrecer apoyo a la persona afectada por la enfermedad.\\
  	
  	\par
  	
  	Existe una versión gratuita, que permite crear una notificación y una grabación con el objetivo de probar si esta se ajusta a las necesidades del usuario. En su versión de pago se eliminan las restricciones y permite conectart la aplicación con un \enquote{Android Smart Watch}, añadiendo la posibilidad de realizar acciones cuando el usuario se despierta de noche y así evitar la desorientación. Como mejora de futuro, se está trabajando en añadir una función para saber cuándo sale de casa.\\
  	\par
  	
  	\textbf{Ventajas:}
 \begin{enumerate}[label=-]
	\item Notificaciones totalmente personalizadas.
	\item Se puede grabar cómo reacciona a los mensajes.
	\item Permite bloquear la aplicación con un pin para que el afectado no pueda configurarla.
\end{enumerate}

	\textbf{Inconvenientes:}
\begin{enumerate}[label=-]
	\item Hay que configurar la aplicación y ver las reacciones desde el dispositivo de la persona que la va a utilizar.
	\item Para funcionalidades más complejas se necesita de un dispositivo Smart Watch.
	\item No se puede personalizar la aplicación.
	\item En caso de olvidar el pin no se puede acceder a la aplicación.
\end{enumerate}

\subsection*{Timeless}

Esta aplicación, creada por Emma Yang en 2018 con tan sólo 14 años, nació con el objetivo dede ayudar a su abuela. Ofrece una conexión en tiempo real entre paciente, cuidador y familia/amigos. Permite enviar fotos con etiquetas de identificación, programar eventos y, con ayuda del localizador del dispositivo, saber el tiempo, una galería de fotos de personas, la opción de identificar automáticamente a la persona a la que se le hace la fotografía desde la app y realizar llamadas a las personas que así se haya configurado.\\

\par

Desde su creación, la aplicación se ha seguido desarrollando y ya se ofrece en varios idiomas para Mac, iPhone y iPad. La financiación para este proyecto que aún sigue en marcha con planes de mejora como añadir gamificación o el uso de inteligencia artificial para detectar hábitos o anomalías en el comportamiento, vino a través de la plataforma \enquote{Indiegogo}, con un presupuesto inicial de $8.837€$.\\

\par

  	\textbf{Ventajas:}
 \begin{enumerate}[label=-]
	\item El cuidador es el que crea y maneja los eventos del paciente.
	\item Las aplicaciones están sincronizadas.
	\item Permite estar en total contacto con el entorno en una única aplicación.
	\item Utiliza Inteligencia Artificial para el reconocimiento de personas dentro del entorno a partir de una fotografía.
	\item Se accede a la palicación a través de la huella digital o reconocimiento facial.
\end{enumerate}

	\textbf{Inconvenientes:}
\begin{enumerate}[label=-]
	\item Únicamente está disponible para dispositivos Apple.
	\item Requiere un aprendizaje por parte del paciente para el manejo.
	\item No permite utilizar la ubicación para saber dónde se encuentra el paciente.
\end{enumerate}


  \section{Conclusiones}\label{sec.conclusiones arte}
  %Análisis propio (sin referencias bibliográficas) de lo que el autor considera sobre todo lo que ha explicado en el apartado de "Estado del arte". En nuestro caso, considerar qué sistema es mejor que otro y en qué aspectos.
  
  Tras analizar múltiples sistemas de asistencia y aplicaciones disponibles diseñadas especialmente para nuestro público objetivo vemos que hay muy buenas ideas, pero que todas dependen de un dispositivo móvil o \enquote{\textit{wearable}}. En el caso de la aplicación \enquote{Timeless} vemos resuelto ese gran problema de la comunicación simple a través de una plicación, en la que los cuidadores manejan casi por completo toda la aplicación del paciente que sufre EA; por otro lado, Jiobit ofrece una localización en tiempo real con un dispositivo que pasa desapercibido y Alexa se puede utilizar para establecer rutinas. Cada uno de los sitemas y aplicaciones vistos están diseñados con un objetivo específico que cumplen con creces, pero ninguno de ellos hace un uso unificado de todas las tecnologías, lo que podría solventar muchas deficiencias que tienen cada uno de ellos con respecto a los otros para solucionar nuestro problema. %incluimos el archivo donde se encuentra el capítulo de estado del arte por puntos

\chapter{Diseño y descripción del sistema}\label{cap.diseno y descripcion del sistema} %En este apartado se detallarían decisiones de diseño e implementación.

\input{capitulos/diseno/requisitos previos}

\input{capitulos/diseno/mockups}


  \section{Conclusiones}\label{sec.conclusiones diseno}%incluimos el archivo donde se encuentra el capítulo de diseño y descripción del sistema por puntos

\chapter{Prototipos y desarrollo}\label{cap.prototipos y desarrollo}
%Poner tantos prototipos como estimes necesario. Si mejoramos algo, crear un prototipo nuevo.
  
\section{Prototipo 1}\label{sec.prototipo1}

\section{Conclusiones}\label{sec.conclusiones prototipo}
%incluimos el archivo donde se encuentra el capítulo de prototipos y desarrollo por puntos

\chapter{Conclusiones y mejoras futuras}\label{cap.conclusiones y mejoras futuras}

  \section{Conclusiones técnicas}\label{sec.conclusiones tecnicas}
  %Hablar desde el punto de vista del software

  \section{Conclusiones personales}\label{sec.conclusiones personales}
  %Valoración personal sobre cada una de las cosas vistas durante el proyecto

  \section{Futuras mejoras}\label{sec.futuras mejoras}
  %Trabajos futuros
  %Acceso a los altavoces del paciente para poder reproducir audios desde su dispositivo que puedan clamarle.
  %pegatinas personalizables totalmente (con ayuda de algún framework o api editor visualmente)
  %Para la app paciente incluir ejercicios diarios de memoria, relación, tareas diarias como recordatorio de pastillas etc%incluimos el archivo donde se encuentra el capítulo de conclusiones finales en español

\chapter{Conclusions and future works}\label{cap.conclusions and future works}
 
 
   \section{Technical conclusions}\label{sec.technical conclusions}

 
   \section{Personal conclusions}\label{sec.personal conclusions}
 
   \section{Future works}\label{sec.future works}%incluimos el archivo donde se encuentra el capítulo de conclusiones finales en español

%usar bibtex (pdflatex, bibtex, pdflatex, pdflatex)
\cleardoublepage
%\cite{Stavropoulos2017} %Prueba de funcionamiento de la bibliografia
\addcontentsline{toc}{chapter}{Bibliografía}
\bibliographystyle{acm} % estilo de la bibliografía.
\bibliography{citas.bib} % citas.bib es el fichero donde está salvada la bibliografía.

\end{document}
